\subsection{非平滑接触与多体时间步进法}
经典刚体动力学领域中，关于约束与接触的处理大致可分为罚函数法（Penalty Method / SMC）和互补法（NSC）。Moreau 和 Jean 提出的基于 LCP 的速度级时间步进法因能保障严格非穿透性被广泛植入至开源与商业引擎（如 Project Chrono, PhysX 等）。在标准步进算法的设计中，碰撞被视作发生在“一阶无限薄平面的切平面空间”，其通过 Baumgarte 稳定化参数 $\alpha$ 按反比 $\Delta t$ 纠正现行误差。多位学者指出此刚性修正极易引发数值爆炸，因而提出了各类基于阻尼衰减或CFM（Constraint Force Mixing）的柔性缓冲。

\subsection{隐式几何与 SDF 碰撞代理}
以 OpenVDB/NanoVDB 为先驱的稀疏体素网格结构极大加快了距离场求解效率，成为现今渲染渲染与碰撞算法的主流基础。研究者常使用球追踪（Sphere Tracing）或单点刺探（Point-probe）对 SDF 进行快速深度评估。然而，在离散网格向亚网格级别内插获取法向量时，常常因跨越体素带来梯度的不连续性。现有的平滑方法多集中在预处理阶段，使用几何平滑滤波器对 SDF 整体网格实行拉普拉斯光滑衰减，但这会丧失物理本体的精确尖锐特征（Sharp features）。

\subsection{对高阶与曲率的先验探索}
在图形学研究中，已有学者注意到了一阶截断的局限性并引入曲率补偿模型，如通过特征映射（Feature matching）预测未来时刻的穿透状态。但这些方法多用于位置级动力学（PBD）的投影视角，在正统以物理量纲标定的脉冲冲量 LCP 框架下，结合连续距离场本征二阶偏导矩阵（Hessian Tensor）作为动位移反馈修正的研究尚属真空。本文致力于填补这一鸿沟。